\section{Analysis on Downstream Applications}
% Here we conduct analysis on two downstream tasks and corresponding models that claim to incorporate commonsense knowledge. 
% Both of the tasks are generative and we evaluate by generating a set of sentences and quantify representational harms using definitions in Section~\ref{sec:formalization}
\subsection{CSKB Completion}
As a popular downstream application, we first consider the task of \textit{commonsense knowledge base completion} which looks to automatically augment a CSKB with generated facts~\cite{li2016commonsense}. We focus our analysis on the COMeT model~\cite{bosselut2019comet}, built by fine-tuning a pre-trained GPT model~\cite{radford2018improving} over ConceptNet triples.
% uses a transformer-based~\cite{vaswani2017attention} language model GPT~\cite{radford2018improving} and is further trained on triples from ConceptNet. 
COMeT has been shown to generate unseen commonsense knowledge in ConceptNet with high quality, and much recent work has used it to provide commonsense background knowledge~\cite{shwartz2020unsupervised,chakrabarty-etal-2020-r}. 

% \subsubsection{COMeT Results}
\para{Data}
We collect statements in COMeT as follows: we input the same target words used in ConceptNet as prompts and collect triples by following all relations existing in the model. Specifically, we collect the top 10 generated results from beam search for all 34 relations existing in COMeT learned from ConceptNet. We generate triples for all the targets we consider, resulting in 112k statements converted from triples and masked target words, the same process as we do for ConceptNet.

\para{Overgeneralization}
From the results of the analysis on statements generated by COMeT, one can observe that the \textit{overgeneralization} issue still exists in the generated statements. For instance for the ``\emph{Religion}'' category, the mean of the negative regard is approximately 25\%. This illustrates the prejudice toward the targets in the religion category in terms of \textit{overgeneralization}. In addition, sentiment scores as high as 50\% for some of the targets in some categories represent the severity of overgeneralization bias. Some additional qualitative examples are also included in Table~\ref{motivation_table_Comet}.

\begin{table}[t]
\centering
\scalebox{0.85}{
\centering
\begin{tabular}{p{1.8cm}|p{6.0cm}}
 \toprule
\textbf{Source}&\textbf{Examples}\\[0.5pt]
 \midrule
  \parbox[t]{2mm}{\multirow{4}{*}{\shortstack[c]{COMeT}}}
  &(mexican, Causes, poverty)\\[0.5pt]
 &(muslim, Causes, terrorism)\\[0.5pt]
 &(policeman, Causes, death)\\[0.5pt]
 &(brother, CapableOf, be mother fxxker)\\[0.5pt]
  \bottomrule
\end{tabular}
}
\caption{Qualitative examples of existing biases in a downstream knowledge generation model COMeT. We can observe how destructive biases also exist in these models. This model should not be generating biased commonsense knowledge (prejudices) towards targets, such as mexican.}
\label{motivation_table_Comet}
\end{table}
\para{Disparity in Overgeneralization}
Notice that in COMeT we do not have the data imbalance problem since COMeT is a generative model, and we generate an equal number of statements for each target. Disparity in number of triples is not an issue for this task. However, the disparity in overgeneralization is still an issue in COMeT. For instance, the results from COMeT shown in Figure \ref{fig:testtest} demonstrate the fact that variances exist in both regard and sentiment measures which is an indication of disparity in overgeneralization. This means that some targets are still extremely favored or disfavored according to regard and sentiment percentages compared to other targets, and that this disparity is still apparent amongst the targets.

\subsection{Neural Story Generation}

As our second downstream task, we consider Commonsense Story Generation (CSG)~\cite{guan2020knowledge}: given a prompt, the model will generate 3 to 5 sentences to tell a story. The CSG model augments GPT-2~\cite{radford2019language} with external commonsense knowledge by training on the CSKB examples constructed from ConceptNet and ATOMIC~\cite{sap2019atomic}. 

\para{Data}
To analyze bias in the story output for CSG, we prompt the CSG model using sentences that are about the social perception of a certain target. We split our targets into: \textit{people}, \textit{locations}, \textit{professions}, and \textit{others}.
Next, we manually come up with 30 templates inspired by the prefix templates for bias in NLG~\cite{sheng2019woman}. Some examples are listed in Table~\ref{tab: CSG_prompts}. We then generate prompts by filling the corresponding templates with target names, resulting in around 3k prompts for CSG. 
% For CSG story generation, we separate generated stories by sentence and treat each sentence individually because sentiment and regard has been proven to work on sentences but not long stories~\cite{sheng2019woman}. 
CSG generates a total of 12k sentences and we calculate regard and sentiment percentages based on all the sentences for a given story.
%\xiang{couple important implementation details can be clarified more; add pointers to appendix?}
\begin{table}[t]
\centering
\scalebox{0.7}{
\begin{tabular}{|c|c|}
\hline
\textbf{Category}                     & \textbf{Template}                                                         \\ \hline
\multirow{2}{*}{People}      & \emph{XYZ} was asked about what did they work as.                       \\ \cline{2-2} 
                             & People are discussing about \emph{XYZ}.                                 \\ \hline
\multirow{2}{*}{Locations}   & People are discussing about people from \emph{XYZ}.                     \\ \cline{2-2} 
                             & People are discussing about traveling to \emph{XYZ}. \\ \hline
\multirow{2}{*}{Professions} & People are talking about what they think of \emph{XYZ}.                 \\ \cline{2-2} 
                             & The \emph{XYZ} was asked about what was they regarded as                \\ \hline
\multirow{2}{*}{Others}      & People are discussing about \emph{XYZ}.                                 \\ \cline{2-2} 
                             & People are talking about what they think of \emph{XYZ}.                 \\ \hline
\end{tabular}
}
\caption{Example prompt templates for story generation for different targets inspired by~\cite{sheng2019woman}. XYZ is replaced with the target name.}
\label{tab: CSG_prompts}
\end{table}

\para{Overgeneralization}
From Figure~\ref{fig:testtest}, we observe similar patterns in terms of the existence of the overgeneralization issue. For instance, as shown in the results in Figure~\ref{fig:testtest}, categories like religion span up to having 60\% negative associations in terms of regard and sentiment scores. 
% This is the case when for ConceptNet these results are more moderate. 
% Thus, we can observe how these types of overgeneralization issues not only exist but are amplified in this downstream task.

\para{Disparity in Overgeneralization}
%\xiang{show metrics by Eq (4) too?}
Similar to the COMeT model since we generated equal amount of statements for this task, we do not observe the disparity in the number of statements as we did with ConceptNet. However, as illustrated in the results presented in Figure~\ref{fig:testtest}, the disparity in overgeneralization is still problematic. For instance, as in Figure~\ref{fig:testtest} the disparity in the ``\emph{Religion}'' category on the negative sentiment spans from 0\% to 60\%. In addition, the ``\emph{Origin}'' category for the CSG task has a significant spread similar to other categories, such as ``\emph{Religion}'' and ``\emph{Gender}''. 


% \begin{table*}
% %\centering
% \begin{tabular}{ p{1.5cm} p{2.0cm} p{1.4cm} p{1.5cm} p{1.5cm} p{1.6cm} p{1.6cm} p{1.6cm}}
%  \toprule
% \textbf{Category}& \textbf{Measure}& \textbf{Negative Mean} & \textbf{Negative Variance} &\textbf{Positive Mean}&\textbf{Positive Variance}& \textbf{Neutral Mean}& \textbf{Neutral Variance}\\[0.5pt]
%  \midrule
%  \parbox[t]{2mm}{\multirow{2}{*}{\shortstack[l]{Race}}}
%  & Sentiment&5.0&75.0&40.0&133.3&54.9&408.3\\[0.5pt]
%  &Regard&5.0&75.0&21.7&186.1&73.3&244.4\\[0.5pt]
%  \midrule
%  \parbox[t]{2mm}{\multirow{2}{*}{\shortstack[l]{Gender}}} & 
%  Sentiment&9.6&134.6&44.3&68.0&46.1&189.0\\[0.5pt]
%  &Regard&26.9&160.3&6.1&119.5&66.9&264.6\\[0.5pt]
%   \midrule
%  \parbox[t]{2mm}{\multirow{2}{*}{\shortstack[l]{Profession}}} &
%  Sentiment&0.9&71.2&52.2&339.5&46.8&358.3\\[0.5pt]
%  &Regard&11.2&314.9&17.3&421.6&71.5&428.3\\[0.5pt]
%  \bottomrule
% \end{tabular}
% \caption{Results from KG-BART.}
% \label{kg_bart_table}
% \end{table*}
\subsection{Bias Mitigation on CSKB Completion}
To mitigate the observed representational harms in ConceptNet and their effects on downstream tasks, we propose a pre-processing data filtering technique that reduces the effect of existing representational harms in ConceptNet. We apply our mitigation technique on COMeT as a case study. 
% However, since our method is a pre-processing method that only augments the data, it can be applied to any other task and model.

\para{Mitigation Approach}
\begin{table}[t]
\centering
\scalebox{0.6}{
\begin{tabular}{ p{2.4cm} cccccc}
 \toprule
\textbf{}&\textbf{NSM $\uparrow$}&\textbf{NSV $\downarrow$}&\textbf{NRM $\uparrow$}&\textbf{NRV $\downarrow$}&\textbf{HNM $\uparrow$}&\textbf{Quality $\uparrow$}\\
 \midrule
 \parbox[t]{2mm}{\multirow{1}{*}{\shortstack[l]{COMeT}}}
&62.6&33.4&78.2&63.2&55.8&\textbf{55.8}\\[0.5pt]

 \parbox[t]{2mm}{\multirow{1}{*}{\shortstack[l]{COMeT-Filtered}}} 
&\textbf{63.2}&\textbf{32.3}&\textbf{78.6}&\textbf{56.7}&\textbf{60.5}&49.9\\[0.5pt]


%  \midrule
 \bottomrule
\end{tabular}
}
\caption{Mitigation results of the filtering technique (COMeT-Filtered) compared to standard COMeT. COMeT-Filtered is effective at reducing overgeneralization and disparity according to sentiment and regard measures and human evaluation. The quality of the generated triples from COMeT, however, is compromised. 
}
\label{sentiment_table}
\end{table}
% \smallskip
% \noindent
% \textbf{Approach~}
Our pre-processing technique relies on data filtering. In this approach, the ConceptNet triples are first passed through regard and sentiment classifiers and only get included in the training process of the downstream tasks if they do not contain representational harms in terms of our regard and sentiment measures. In other words, in this framework, all the biased triples that were associated with a positive or negative label from regard and sentiment classifiers get filtered out and only neutral triples with neutral label get used. 
% We then filter this data according to our technique and report the results in Table \ref{sentiment_table}. 
% \begin{table*}[t]
% \centering
% \scalebox{0.99}{
% \begin{tabular}{ p{2.7cm} p{1.5cm} p{1.5cm} p{1.5cm} p{1.5cm} p{1.5cm} p{2.0cm}}
%  \toprule
% \textbf{Model}& \textbf{NSM $\uparrow$} &\textbf{NSV $\downarrow$}&\textbf{NRM $\uparrow$}&\textbf{NRV $\downarrow$}&\textbf{HNM $\uparrow$}&\textbf{Quality $\uparrow$}\\[0.5pt]
%  \midrule
%  \parbox[t]{2mm}{\multirow{1}{*}{\shortstack[l]{COMeT}}}
%  &62.6&33.4&78.2&63.2&55.8&\textbf{55.8}\\[0.5pt]
% %  \midrule
%  \parbox[t]{2mm}{\multirow{1}{*}{\shortstack[l]{COMeT-Filtered}}} 
%  &\textbf{63.2}&\textbf{32.2}&\textbf{78.6}&\textbf{56.7}&\textbf{60.5} & 49.9\\[0.5pt]
% %  \midrule
%  \bottomrule
% \end{tabular}
% }
% \caption{Mitigation results for filtering technique (COMeT-Filtered) compared to standard COMeT. We find the simple filtering approach is effective at reducing both overgeneralization and disparity according to our sentiment and regard measures as well as human evaluation. The quality of the generated triples from COMeT, however, is compromised.}
% \label{sentiment_table}
% \end{table*}

% \begin{table*}[t]
% \centering
% \scalebox{0.99}{
% \begin{tabular}{ p{2.2cm} p{1.5cm} p{2cm} p{3cm} p{3cm}}
%  \toprule
% \textbf{}&\textbf{COMeT Regard}& \textbf{COMeT Sentiment} &\textbf{COMeT-Filtered Regard}&\textbf{COMeT-Filtered Sentiment}\\[0.5pt]
%  \midrule
%  \parbox[t]{2mm}{\multirow{1}{*}{\shortstack[l]{\textbf{Human Label}}}}
%  &\textcolor{red}{71.6\%}&\textcolor{red}{59.2\%}&\textcolor{red}{72.8\%}&\textcolor{red}{54.7\%}\\[0.5pt]
% %  \midrule

% %  \midrule
%  \bottomrule
% \end{tabular}
% }
% \caption{Percentages represent how much regard and sentiment labels ran on COMeT and COMeT-Filtered triples agree with labels coming from humans. The higher the percentage, it means that the measure agrees with human's perception of bias more closely. The results show that regard might be better aligned with human's perception of bias and can serve as a good proxy to measure biases in triples coming from both models.}
% \label{human_vs_re_sent}
% \end{table*}

% \begin{table}[t]
% \centering
% \scalebox{0.94}{
% \begin{tabular}{ p{1.4cm} cc}
%  \toprule
% \textbf{}&\textbf{COMeT} &\textbf{COMeT-Filtered}\\
%  \midrule
%  \parbox[t]{2mm}{\multirow{1}{*}{\shortstack[l]{NSM $\uparrow$}}}
% &62.6&\textbf{63.2}\\[0.5pt]

%  \parbox[t]{2mm}{\multirow{1}{*}{\shortstack[l]{NSV $\downarrow$}}} 
% &33.4&\textbf{32.2}\\[0.5pt]

%  \parbox[t]{2mm}{\multirow{1}{*}{\shortstack[l]{NRM $\uparrow$}}}
%  &78.2&\textbf{78.6}\\[0.5pt]

%  \parbox[t]{2mm}{\multirow{1}{*}{\shortstack[l]{NRV $\downarrow$}}}
% &63.2&\textbf{56.7}\\[0.5pt]

%  \parbox[t]{2mm}{\multirow{1}{*}{\shortstack[l]{HNM $\uparrow$}}}
%  &55.8&\textbf{60.5}\\[0.5pt]

%  \parbox[t]{2mm}{\multirow{1}{*}{\shortstack[l]{Quality $\uparrow$}}}
%  &\textbf{55.8}&49.9\\[0.5pt]
% %  \midrule
%  \bottomrule
% \end{tabular}
% }
% \caption{Mitigation results for filtering technique (COMeT-Filtered) compared to standard COMeT. We find the simple filtering approach is effective at reducing both overgeneralization and disparity according to our sentiment and regard measures as well as human evaluation. The quality of the generated triples from COMeT, however, is compromised. 
% \xiang{can rotate the table to save space.}
% }
% \label{sentiment_table}
% \end{table}


% \begin{table}[t]
% \centering
% \scalebox{0.8}{
% \begin{tabular}{ p{1.5cm} C{1.5cm} C{3.0cm}}
%  \toprule
% \textbf{Metrics}& \textbf{COMeT} &\textbf{COMeT-Filtered}\\[0.5pt]
%  \midrule
%  \parbox[t]{2mm}{\multirow{1}{*}{\shortstack[l]{NSM $\uparrow$}}}
%  &62.6&\textbf{63.2} \\[0.5pt]
% %  \midrule
%  \parbox[t]{2mm}{\multirow{1}{*}{\shortstack[l]{NSV $\downarrow$}}} &
%  33.4&\textbf{32.2}\\[0.5pt]
% %  \midrule
%  \parbox[t]{2mm}{\multirow{1}{*}{\shortstack[l]{NRM $\uparrow$}}} & 
% 78.2&\textbf{78.6}\\[0.5pt]
% %   \midrule
%  \parbox[t]{2mm}{\multirow{1}{*}{\shortstack[l]{NRV $\downarrow$}}} &
% 63.2&\textbf{56.7}\\[0.5pt]
%  \bottomrule
% \end{tabular}
% }
% \caption{Mitigation results for filtering technique compared to vanilla COMeT.}
% \label{sentiment_table}
% \end{table}


% \begin{table}
% %\centering
% \begin{tabular}{ p{2.2cm} C{1.5cm} C{2.8cm}}
%  \toprule
% \textbf{Measure}& \textbf{COMeT} & \textbf{COMeT-Filtered}\\[0.5pt]
%  \midrule
%  \parbox[t]{2mm}{\multirow{1}{*}{\shortstack[l]{Neutral Mean $\uparrow$}}}
%  &55.8&\textbf{60.5}\\[0.5pt]
%  \midrule
%  \parbox[t]{2mm}{\multirow{1}{*}{\shortstack[l]{Quality $\uparrow$}}}
%  & \textbf{55.8}&49.9\\[0.5pt]
%  \bottomrule
% \end{tabular}
% \caption{Human annotator results.}
% \label{human_detail_table}
% \end{table}

% \subsection{Data Augmentation}
% In our second pre-processing mitigation technique, instead of filtering triples, we include all the triples and level off all the groups with non-neutral triples. This method can complement the previous method in which it does not remove triples; thus, still concepts are present and no information is lost. In this approach, triples that are detected to be positive or negative for a certain demographic are replicated for all the other demographics in the category that the detected negative/positive triple appeared in so that all the demographics are leveled off in that category. For instance, for the triple $(\text{American, IsA, great\_person})$ in the race category after its conversion to ``\emph{American is a great person}'', we pass it through the sentiment and regard classifiers. If this sentence is detected to be un-neutral, we replicate this triple to other existing demographics in the race category, such as ``\emph{Chinese}', and create the new instances ``\emph{Chinese is a great person}'' same as for other demographics in the race category and add it to the augmented data. This augmented data with leveled off demographics in a certain category is then used to train downstream tasks.

% \smallskip
% \noindent
% \textbf{Results}
% \subsubsection{Results on Bias Metrics}
% \vspace{-0.1cm}

% To analyze the results of our mitigation technique, we consider both overgeneralization and disparity in overgeneralizatio over which we reduce the existing biases. 

\para{Results on Overgeneralization} To measure effectiveness of mitigation over overgeneralization, we consider increasing the overall mean of neutral triples which is indicative of reducing the overall favoritism and prejudice according to sentiment and regard measures. We report the effects on overgenaralization on sentiment as \textit{Neutral Sentiment Mean (NSM)} and regard measure as \textit{Neutral Regard Mean (NRM)}. As demonstrated in Table \ref{sentiment_table}, by increasing the overall neutral sentiment and regard means, our filtered model is able to reduce the unwanted positive and negative associations and reduce the overgeneralization issue.

\para{Results on Disparity in Overgeneralization} To measure effectiveness of mitigation over disparity in overgeneralization, we consider reducing the existing variance amongst different targets. We report the disparity in overgeneralization on sentiment as \textit{Neutral Sentiment Variance (NSV)} and on regard as \textit{Neutral Regard Variance (NRV)}. Shown in Table \ref{sentiment_table}, our filtered technique reduces the variance and disparities amongst targets over the standard COMeT model in terms of regard and sentiment measures.  

% However, we acknowledge the fact that these improvements are minor; thus, we encourage future work to focus on improvement of these results and propose stronger mitigation techniques.  

%\vspace{-0.1cm}
\para{Human Evaluation of Mitigation Results}\label{human_eval}
%\vspace{-0.1cm}
% \para{Human Evaluation} 
In addition to reporting regard and sentiment scores, we perform human evaluation on 3,000 generated triples from standard COMeT and COMeT-Filtered  models to evaluate both the quality of the generated triples and the bias aspect of it from the human perspective on Amazon Mechanical Turk. From the results in Table \ref{sentiment_table}, one can observe that COMeT-Filtered is construed to have less overall overgeneralization harm since humans rated more of the triples generated by it to be neutral and not containing negative or positive associations. This is shown as \textit{Human Neutral Mean (HNM)} in Table~\ref{sentiment_table}. However, this came with a trade-off for quality in which COMeT-Filtered is rated to have less quality compared to standard COMeT in terms of validity of its triples. We encourage future work to improve for higher quality. In addition, we measure the inter-annotator agreement and report the Fleiss' kappa scores \cite{fleiss1971measuring} to be 0.4788 and 0.6407 for quality and representational harm ratings respectively in the standard COMeT model and 0.4983 and 0.6498 for that of COMeT-Filtered.